
\documentclass{article}
\usepackage{colortbl}
\usepackage{makecell}
\usepackage{multirow}
\usepackage{supertabular}

\begin{document}

\newcounter{utterance}

\centering \large Interaction Transcript for game `cladder', experiment `full\_v1.5\_default', episode 9560 with refineonly{-}Qwen3.5{-}4B.
\vspace{24pt}

{ \footnotesize  \setcounter{utterance}{1}
\setlength{\tabcolsep}{0pt}
\begin{supertabular}{c@{$\;$}|p{.15\linewidth}@{}p{.15\linewidth}p{.15\linewidth}p{.15\linewidth}p{.15\linewidth}p{.15\linewidth}}
    \# & \multicolumn{2}{c}{Player} && \multicolumn{2}{c}{Game Master} \\
    \hline

    \theutterance \stepcounter{utterance}  
    & & & \multicolumn{4}{p{0.6\linewidth}}{
        \cellcolor[rgb]{0.9,0.9,0.9}{
            \makecell[{{p{\linewidth}}}]{
                \texttt{\tiny{[P1$\langle$GM]}}
                \texttt{You are an expert in causal inference. The following question is not a typical commonsense query, but rather a meticulously designed question created by a professor specializing in causal inference, intended to assess the students' mastery of the course content.} \\
\texttt{Imagine a self{-}contained, hypothetical world with only the following conditions, and without any unmentioned factors or causal relationships: Tijv has a direct effect on xevo and gyzp. Xevo has a direct effect on gyzp. We know that tijv causes not xevo. tijv and xevo causes gyzp. We observed an individual is not tijv. Would an individual is gyzp if not xevo instead of xevo?} \\
\texttt{Start your answer with 'Yes' or 'No', followed by additional reasoning or evidence to support your explanation.} \\
            }
        }
    }
     \\ \\

    \theutterance \stepcounter{utterance}  
    & & \multicolumn{4}{p{0.6\linewidth}}{
        \cellcolor[rgb]{0.95,0.95,0.95}{
            \makecell[{{p{\linewidth}}}]{
                \texttt{\tiny{[P1$\leftrightarrow$P]}}
                \texttt{{[}DRAFT{]}} \\
\texttt{No} \\
\\ 
\texttt{**Reasoning:**} \\
\\ 
\texttt{To determine the answer, we must construct the causal structure based strictly on the provided conditions and apply the rules of causal inference.} \\
\\ 
\texttt{1.  **Construct the Causal Graph:**} \\
\texttt{    *   **Tijv** has a direct effect on **xevo**.} \\
\texttt{    *   **Tijv** has a direct effect on **gyzp**.} \\
\texttt{    *   **Xevo** has a direct effect on **gyzp**.} \\
\texttt{    *   *Note on Contradiction Resolution:* The prompt states "Tijv causes not xevo" and "Tijv and xevo causes gyzp." In causal diagrams, "Tijv causes not xevo" implies a negative edge (\$Tijv \textbackslash{}to \textbackslash{}neg xevo\$). However, the first sentence says "Tijv has a direct effect on xevo." In standard causal notation, a direct effect is usually positive unless specified as negative. Given the explicit statement "Tijv causes not xevo," we interpret the relationship as \$Tijv \textbackslash{}to \textbackslash{}neg xevo\$.} \\
\texttt{    *   The statement "Tijv and xevo causes gyzp" implies that both variables are required (or contribute) to the outcome. This creates a structure where \$gyzp\$ is a function of \$Tijv\$ and \$xevo\$.} \\
\\ 
\texttt{2.  **Analyze the Counterfactual Query:**} \\
\texttt{    *   **Observed State:** The individual is **not tijv** (\$Tijv = 0\$).} \\
\texttt{    *   **Question:** Would the individual be **gyzp** if **not xevo** (\$xevo = 0\$) instead of **xevo** (\$xevo = 1\$)?} \\
\texttt{    *   This is a counterfactual question: \$P(gyzp=1 | do(xevo=0), Tijv=0)\$.} \\
\\ 
\texttt{3.  **Evaluate the Causal Mechanism:**} \\
\texttt{    *   Since \$Tijv = 0\$, the negative effect on \$xevo\$ is nullified (or \$xevo\$ is determined by other factors not mentioned, but we must assume the system is closed as per the prompt).} \\
\texttt{    *   The outcome \$gyzp\$ is caused by \$Tijv\$ and \$xevo\$.} \\
\texttt{    *   If \$Tijv = 0\$, the contribution of the first cause to \$gyzp\$ is absent (assuming a positive contribution from the variable itself, which is standard unless "not" is part of the variable name).} \\
\texttt{    *   The query asks what happens if we set \$xevo = 0\$ (instead of 1).} \\
\texttt{    *   If \$gyzp\$ requires *both* \$Tijv\$ and \$xevo\$ to be present (as implied by "Tijv and xevo causes gyzp"), then setting \$} \\
            }
        }
    }
    & \\ \\

    \theutterance \stepcounter{utterance}  
    & \multicolumn{4}{p{0.6\linewidth}}{
        \cellcolor[rgb]{0.561,0.737,0.561}{
            \makecell[{{p{\linewidth}}}]{
                \texttt{\tiny{[P1$\rangle$GM]}}
                \texttt{No} \\
\\ 
\texttt{**Reasoning:**} \\
\\ 
\texttt{To determine the answer, we must construct the causal structure based strictly on the provided conditions and apply the rules of causal inference.} \\
\\ 
\texttt{1.  **Construct the Causal Graph:**} \\
\texttt{    *   **Tijv** has a direct effect on **xevo**.} \\
\texttt{    *   **Tijv** has a direct effect on **gyzp**.} \\
\texttt{    *   **Xevo** has a direct effect on **gyzp**.} \\
\texttt{    *   *Note on Contradiction Resolution:* The prompt states "Tijv causes not xevo" and "Tijv and xevo causes gyzp." In causal diagrams, "Tijv causes not xevo" implies a negative edge (\$Tijv \textbackslash{}to \textbackslash{}neg xevo\$). However, the first sentence says "Tijv has a direct effect on xevo." In standard causal notation, a direct effect is usually positive unless specified as negative. Given the explicit statement "Tijv causes not xevo," we interpret the relationship as \$Tijv \textbackslash{}to \textbackslash{}neg xevo\$.} \\
\texttt{    *   The statement "Tijv and xevo causes gyzp" implies that both variables are required (or contribute) to the outcome. This creates a structure where \$gyzp\$ is a function of \$Tijv\$ and \$xevo\$.} \\
\\ 
\texttt{2.  **Analyze the Counterfactual Query:**} \\
\texttt{    *   **Observed State:** The individual is **not tijv** (\$Tijv = 0\$).} \\
\texttt{    *   **Question:** Would the individual be **gyzp** if **not xevo** (\$xevo = 0\$) instead of **xevo** (\$xevo = 1\$)?} \\
\texttt{    *   This is a counterfactual question: \$P(gyzp=1 | do(xevo=0), Tijv=0)\$.} \\
\\ 
\texttt{3.  **Evaluate the Causal Mechanism:**} \\
\texttt{    *   Since \$Tijv = 0\$, the negative effect on \$xevo\$ is nullified (or \$xevo\$ is determined by other factors not mentioned, but we must assume the system is closed as per the prompt).} \\
\texttt{    *   The outcome \$gyzp\$ is caused by \$Tijv\$ and \$xevo\$.} \\
\texttt{    *   If \$Tijv = 0\$, the contribution of the first cause to \$gyzp\$ is absent (assuming a positive contribution from the variable itself, which is standard unless "not" is part of the variable name).} \\
\texttt{    *   The query asks what happens if we set \$xevo = 0\$ (instead of 1).} \\
\texttt{    *   If \$gyzp\$ requires *both* \$Tijv\$ and \$xevo\$ to be present (as implied by "Tijv and xevo causes gyzp"), then setting \$} \\
            }
        }
    }
    & & \\ \\

    \theutterance \stepcounter{utterance}  
    & & & \multicolumn{2}{p{0.3\linewidth}}{
        \cellcolor[rgb]{0.9,0.9,0.9}{
            \makecell[{{p{\linewidth}}}]{
                \texttt{\tiny{[GM$|$GM]}}
                \texttt{No} \\
            }
        }
    }
    & & \\ \\

    \theutterance \stepcounter{utterance}  
    & & & \multicolumn{2}{p{0.3\linewidth}}{
        \cellcolor[rgb]{0.9,0.9,0.9}{
            \makecell[{{p{\linewidth}}}]{
                \texttt{\tiny{[GM$|$GM]}}
                \texttt{no} \\
            }
        }
    }
    & & \\ \\

    \theutterance \stepcounter{utterance}  
    & & & \multicolumn{2}{p{0.3\linewidth}}{
        \cellcolor[rgb]{0.9,0.9,0.9}{
            \makecell[{{p{\linewidth}}}]{
                \texttt{\tiny{[GM$|$GM]}}
                \texttt{game\_result = WIN} \\
            }
        }
    }
    & & \\ \\

\end{supertabular}
}

\end{document}
